\section{Flat-domain Gaussian limits}
\label{sec:gaussian}

We prove the Gaussian limit for the unit-increment height
\(K_\delta=2H_\delta\), first in deterministic flat chambers and then
conditionally in chambers produced by a stopping exploration. The uniformity
in the chamber is what permits the second use.

Every chamber in this section is an accessible component behind a connected
matching \(W\) perimeter, with the medial resolution at pinches from
\cref{lem:finite-cuts}. Its open face domain is simply connected. The
perimeter carries its actual wire, and incident faces are not joined across a
pinch. We subtract its constant height. For the moment arguments choose \(N\)
to be the free current complementary to this wall percolation \(W\); all its
interior components then have fair colors. The initial flat domain is
included by taking \(N\) to be the free dual current. Either corner parity
may carry the wall.

A collection of chambers will always have pairwise disjoint open interiors;
their closures may touch. Sample their centered flat laws independently,
patch the fields on these disjoint cell interiors, and set the field to zero
on the exposed complement. Let \(U_\delta\) denote their open union. Fix an
open disk \(B_0\) containing all chamber closures with a fixed outer margin,
and put
\[
             F_\delta=\overline B_0\setminus U_\delta.
\]
Thus \(F_\delta\) is the closed complement in a common compact ambient set.
Hausdorff convergence of these \(F_\delta\) is equivalent here to local
Hausdorff convergence of the closed complements in the plane. Constants may
depend on \(B_0\), the fixed parameter, and the moment order, but not on the
shapes or the number of chambers.

\begin{theorem}[Flat-domain limit]\label{thm:flat-gff}
Let \(U_\delta\) be a single deterministic flat chamber, or a finite
mesh-dependent union of independent chambers with pairwise disjoint
interiors, as above. Suppose \(F_\delta\) converges in Hausdorff distance to
\(F\), and set \(U=B_0\setminus F\). Then
\[
                K_\delta\ \Longrightarrow\ h_U/a
             \quad\text{in }H^s_{\mathrm{loc}}(U),\qquad s<-1,
\]
where the restrictions of \(h_U\) to the components of \(U\) are independent
zero-Dirichlet GFFs with logarithmic singularity \(-\log|x-y|\). All finite
joint smooth-test moments with supports in \(U\) converge. Different limiting
components are independent even if their discrete precursors were the same
chamber. Either wall parity and either atomic or interpolated smearing gives
the same limit. The fields extended by zero off \(U_\delta\) are uniformly
tight in \(H^s_{\mathrm{loc}}(B_0)\) for every \(s<-1\).
\end{theorem}

The proof first bounds all moments and their boundary traces. Screened
orthogonality then makes the limiting moments harmonic away from collisions;
a local neutral-pair comparison identifies their logarithmic poles and hence
their Wick recursion. We determine its coefficient from the plane limit
before completing the theorem. The last proposition turns the deterministic
result into the conditional statement needed for stopped explorations.

\subsection{Limiting chambers and their Green kernels}

For an open set \(V\) whose components are simply connected, write \(G_V\)
for the Dirichlet kernel with logarithmic singularity \(-\log|x-y|\) in
each component, extended by zero between different components and when
either argument lies outside \(V\). For tests put
\[
 G_V(f,g)=\iint f(x)G_V(x,y)g(y)\,\dd x\,\dd y.
\]
The planar connected-complement criterion shows that \(F_\delta\) above is
connected: an open subset of the sphere has simply connected components
exactly when its complement is connected. This uses disjoint interiors, not
disjoint closures. The next lemma supplies the boundary facts used twice
below, first in Wick's rule and then in the conditional limit.

\begin{lemma}[Connected complements and Green convergence]
\label{lem:flat-green-geometry}
Let \(F_n\) be connected compact subsets of \(\overline B_0\) containing
\(\partial B_0\), and suppose \(F_n\to F\) in Hausdorff distance. Put
\(V_n=B_0\setminus F_n\) and \(V=B_0\setminus F\). Every component of
\(V_n\) and \(V\) is simply connected, every compact subset of \(V\) is
eventually contained in \(V_n\), and
\[
                G_{V_n}(x,y)\longrightarrow G_V(x,y)
\]
locally uniformly for distinct \(x,y\in V\), including when they lie in
different components of \(V\).
The zero-extended kernels satisfy \(0\le G_{V_n},G_V\le G_{B_0}\).
As a function of the closed complement and its two arguments, this
zero-extended Green kernel is Borel measurable.

For any proper simply connected domain \(W\), with
\(d_W(x)=\dist(x,\partial W)\), one has
\begin{equation}
 0\le G_W(x,y)\le
       4\sqrt{\frac{d_W(x)}{|x-y|}}
 \quad\text{if }d_W(x)\le\frac{|x-y|}{16}.
 \label{eq:flat-green-boundary}
\end{equation}
In particular these Green kernels have zero boundary trace without a
boundary smoothness assumption.
\end{lemma}

\begin{proof}
Adjoin the exterior of \(B_0\) in the sphere. The complements remain
connected, and so does their Hausdorff limit. The planar criterion just
stated gives simple connectivity of every open component. Hausdorff
convergence of compact sets gives uniform convergence of their distance
functions, which proves compact containment.

For the estimate, choose a conformal map \(f:\mathbb D\to W\) with
\(f(0)=x\), and write \(f(z)=y\). The growth bound and the quarter theorem
recalled in \cref{sec:conformal-growth} give
\[
 |x-y|\le |f'(0)|\frac{|z|}{(1-|z|)^2}
          \le \frac{4d_W(x)}{(1-|z|)^2}.
\]
Under the stated hypothesis this implies \(|z|\ge1/2\). The conformal
Green identity and \(-\log t\le2(1-t)\) for \(1/2\le t\le1\) now give
\cref{eq:flat-green-boundary}.

For convergence, domain monotonicity gives the common bound
\(0\le G_{V_n}\le G_{B_0}\), with zero understood between components.
The upper kernel is bounded on compact sets away from the diagonal and
has only an integrable logarithmic singularity on the diagonal. Fix \(y\in V\) and a
component \(C\) of \(V\). On compact subsets of \(C\setminus\{y\}\),
harmonic compactness gives subsequential limits in the first variable.
If \(y\in C\), a fixed disk about \(y\) is eventually contained in one
component of \(V_n\). On that disk subtract \(\log(1/|x-y|)\).
The difference is harmonic and uniformly bounded on a fixed interior
circle by the common upper bound; the Poisson formula therefore preserves
the unit logarithmic pole in the limit. If \(y\notin C\), there is no pole.

For \(x\in V_n\), its distance to the boundary of its component equals
\(\dist(x,F_n)\). Passing \cref{eq:flat-green-boundary} to a subsequential
limit shows that this limit has zero trace on \(\partial C\), since
\(\dist(x,F_n)\to\dist(x,F)\). The maximum principle now identifies it
as \(G_C(\,\cdot\,,y)\) when \(y\in C\), and as zero when \(y\notin C\).
The removable difference is bounded near its pole and near the boundary,
so the maximum principle applies on the whole component. Every
subsequential limit is the same. Interior harmonic estimates in both
variables give the asserted local uniform convergence. The common
logarithmic upper bound also gives convergence of \(G_{V_n}(f,g)\) for
fixed smooth tests with supports compactly contained in \(V\).
On the space of connected complements, the local uniform assertion gives
continuity of \(G_{B_0\setminus F}(x,y)\) where \(x,y\notin F\) and \(x\ne y\).
These positive-clearance sets form a countable cover of the interior
off-diagonal pairs. Extension by zero off the open set and the usual
diagonal value therefore give a Borel kernel, and hence Borel smeared
covariances.
\end{proof}

\subsection{Moments and boundary traces}

Choose once \(R_0>2\diam(B_0)\).

\begin{proposition}[Logarithmic moments]
\label{prop:log-moments}
Let $z_1,\ldots,z_{2m}$ be corner-grid sites in a flat chamber of
diameter at most $R_0/2$. They may lie on either parity. Then
\begin{equation}
 \left|\E\prod_{i=1}^{2m}K(z_i)\right|
 \le C_m\sum_{\mathcal P}\prod_{\{i,j\}\in\mathcal P}
 \left(1+\log^+\frac{R_0}{|z_i-z_j|\vee\delta}\right),
 \label{eq:flat-log-moments}
\end{equation}
where the sum is over pairings. Odd moments vanish. The same bound holds for
independent flat chambers after setting the field to zero off those chambers.
For a compactly supported bounded test $f$, uniformly over all such chambers,
\begin{equation}
 \E|K_\delta(f)|^{2m}\le C_m\|f\|_\infty^{2m}.
 \label{eq:flat-smearing-moments}
\end{equation}
For any fixed smooth compactly supported profile $f$, put
$f_{r,z}(x)=r^{-2}f((x-z)/r)$. If $0<\delta\le r$, then
\begin{equation}
 \E|K_\delta(f_{r,z})|^{2m}
       \le C_{m,f}\left(1+\log^+\frac{R_0}{r}\right)^m.
 \label{eq:flat-scaled-moments}
\end{equation}
This includes a mass-one smooth average of radius $r$ for a fixed profile.
\end{proposition}

\begin{proof}
Condition on the current $N$ and the $s$ spins. By \cref{lem:finite-cuts},
the centered height at a wall-parity site is a sum of jumps of size two with
independent fair signs, one sign per contributing cluster. A cluster's
coefficient at a site is zero or has absolute value two. Its sign is shared
by all its odd holes.

The same signed expansion controls opposite-parity sites. For such a site
$x$, choose an incident wall-parity site $v$, and let $C_x$ be the free
$N$ component containing $x$. The corner rule gives
\[
          K(x)=K(v)+\varepsilon_W s(v)d(C_x),
\]
where the fixed sign $\varepsilon_W$ records which parity carries $W$.
Absorb this term into the coefficient of $d(C_x)$. Conditional on $(N,s)$,
every height is still a sum of the independent fair component colors,
now with coefficients of absolute value at most three. If a component has
a nonzero coefficient at a site, that site lies in the component or its
filled hull up to distance $O(\delta)$: wall-parity contributions come from
odd holes, and the extra contribution is incident to the component.
Consequently a component shared by two sites has diameter at least their
separation minus $O(\delta)$.

Expanding an even product and summing the independent signs leaves only
terms in which every sign appears an even number of times.
Pair the occurrences of each such sign. Overcounting by all pairings gives
an upper bound by $C_m$ times sums of products of shared-cluster counts
$M(z_i,z_j)$. This argument does not assert independence of the counts.

A contributing cluster surrounds or is incident to the relevant site.
The preceding support observation gives the required diameter lower bound
for a shared cluster. At a fixed site only boundedly many components of
diameter at most $C\delta$ can contribute, including isolated vertices;
put all of these contributions in the lattice cutoff. Above that cutoff,
a fixed number of boxes at a cluster's diameter scale detect its passage
or enclosure. By
\cref{prop:cluster-counts}, the count in each such dyadic shell has uniformly
bounded moments of every fixed order. Minkowski's inequality therefore gives
\[
 \|M(z_i,z_j)\|_{L^m}
 \le C_m\left(1+\log^+\frac{R_0}{|z_i-z_j|\vee\delta}\right).
\]
Holder's inequality for the product of the $m$ counts proves
\cref{eq:flat-log-moments}. Height reflection about the flat value proves
vanishing of odd moments. The argument is unchanged in several independent
chambers; shared-cluster counts between different chambers are zero.

Integrate the estimate against the interpolation weights, or the atomic
Riemann weights, of the tests. Each logarithmic factor is uniformly
integrable in its two variables, including its lattice diagonal: the latter
contributes $O(\delta^2\log(1/\delta))$. Since the pairs have disjoint sets
of variables, this proves \cref{eq:flat-smearing-moments}. Rescaling a
fixed profile by $x=z+ru$ gives \cref{eq:flat-scaled-moments}. Here
$\delta/r\le1$ is essential for an atomic smearing: after the rescaling
the diagonal part has the additional uniformly bounded factor
$(\delta/r)^2\log(r/\delta)$, while the common scale contributes
$1+\log^+(R_0/r)$.
\end{proof}

\begin{proposition}[Boundary decay]
\label{prop:boundary-decay}
Consider deterministic flat chambers in a fixed bounded region, with
connected accessible walls, and a Hausdorff limit of their complements.
For any limiting open component, moment kernels obtained from
\cref{eq:flat-log-moments} have zero boundary trace in a distinguished
variable when the other variables remain at distinct interior points.
The bounds and the trace conclusion are uniform for the discrete chambers.
More precisely, a shared cluster joining the height contribution at a site
within distance $r$ of a matching connected wall to one at distance at least
$R$ from that site has probability at most $C(r/R)^\eta$, with fixed
positive outer clearance and $\delta\le r\ll R$.
\end{proposition}

\begin{proof}
The relevant wall component reaches from the small scale to the outside of
the fixed testing window. A shared cluster either is incident to the
distinguished site or encloses it. In the incident case its connected
passage reaches the outer scale forced by the distant spectator. In both cases it must
avoid matching circuits attached to this wall at every scale between
$r$ and $R$: such a circuit separates the distinguished contribution
from the distant one. The conditional circuit estimate in \cref{prop:circuits}
gives a uniformly positive success probability at a fixed fraction of these
separated scales. An attached circuit blocks the cluster, so iteration gives
$C(r/R)^\eta$. The argument is within the accessible chamber with its wire
retained; no regularity of the geometric wall is used.

In every pairing contributing to a mixed even moment, the distinguished
site is paired with one spectator. The corresponding shared cluster has
a fixed positive minimum diameter. Its count, restricted to surviving the
attached-circuit trials, tends to zero in every smaller $L^p$ norm: apply
Holder to the survival indicator and the higher count-moment bound from
\cref{prop:cluster-counts}. The remaining counts are bounded by
\cref{eq:flat-log-moments}. This proves vanishing of the mixed moment as
the distinguished site approaches the wall. For fixed distinct spectators
the pairing bound itself is uniformly bounded near the wall. Smearing first
at a radius smaller than the site's wall distance gives the same result
for distributional moment limits. Once harmonicity is established below,
the mean-value property identifies these smearing limits with the pointwise
boundary trace. The proof also applies when the wall develops a pinch:
only the actual connected wall crossing the testing scales is used.
\end{proof}

\begin{lemma}[Parity and interpolation]
\label{lem:parity-interpolation}
On compact subsets with positive clearance from the wall, the two parity
smearings of $K$, and their smooth nodal interpolation, have the same
subsequential distributional limits and the same joint test moments.
Here each original parity has Riemann weight $\delta^2$; the combined
corner grid has weight $\delta^2/2$ per site. Thus its interpolation is
compared to the average of the two parity smearings, not their sum.
\end{lemma}

\begin{proof}
Pair each site of one parity with an incident site of the other parity by a
fixed lattice translation. The paired height difference is $Z_e\in\{-1,1\}$.
Conditional on $(N,s)$ it is, up to a deterministic orientation, a fixed
$s$ spin times the fair color of the adjacent $N$ component. Thus for two
well-separated paired edges,
\[
 |\E[Z_eZ_{e'}]|\le
 \P(\text{their adjacent sites belong to the same }N\text{ component}).
\]
A matching $W$ circuit around the first edge separates these components.
Bulk circuit trials between the lattice scale and their separation imply
that, for separation at least $r$ and fixed wall clearance,
$|\E[Z_eZ_{e'}]|\le C(\delta/r)^\eta$. This is equivalently the conditional
fair-color coupling in disjoint pin-free buffers. In particular it does not
infer smallness merely from the deterministic bound $|Z_e|=1$.

For a smooth test, the variance of its averaged paired differences is bounded
by $C r^2+C(\delta/r)^\eta$: separate pairs at distance below $r$ from all
other pairs. First let $\delta$ decrease to zero and then let $r$ decrease
to zero. Changing the test's evaluation point within a paired cell gives an
$O(\delta)$ smooth-test error in $L^2$ by
\cref{eq:flat-smearing-moments}. This proves convergence of the difference
to zero in $L^2$. Uniform higher moments and interpolation give convergence
in every finite $L^p$. Nodal interpolation is an average of the same local
pairings with smooth weight errors, so the argument includes it.
\end{proof}

\begin{lemma}[Tightness and absence of collision measures]
\label{lem:flat-tightness}
Centered flat heights are tight in $H^s_{\mathrm{loc}}$, for every $s<-1$,
uniformly over the chambers above. Every finite joint test moment converges
along distributional convergence. Once the limiting moment distributions
are harmonic off their collision diagonals, their smooth representatives
there represent the full distributions: there are no additional measures
or distributions supported on those diagonals.
\end{lemma}

\begin{proof}
Put a cutoff $\chi$ in a fixed square regarded as a torus, with normalized
Fourier basis $e_n$. The uniform second-moment bound for
$K_\delta(\chi e_n)$ is independent of $n$. Hence, for $q>1$,
\[
 \sup_\delta\E\|\chi K_\delta\|_{H^{-q}}^2
 \le C\sum_{n\in\Z^2}(1+|n|^2)^{-q}<\infty.
\]
For $1<q<-s$, compactness of $H^{-q}\hookrightarrow H^s$ gives tightness.
Use a countable exhaustion of cutoffs and summable exceptional probabilities.
Higher versions of \cref{eq:flat-smearing-moments} give uniform
integrability of all finite test products, proving convergence of moments.
These statements also hold for a finite independent collection of chambers or
for fields set to zero on their already exposed complement.

Here is a direct verification of the last assertion, independent of any
pointwise choice on diagonals. Mollify each variable of a $j$th moment
distribution at the same radius $r$. For each fixed $r$ the mesh limit is
taken first, so eventually $\delta\le r$. Holder and
\cref{eq:flat-scaled-moments} bound the mollified kernel by
$C_j(1+|\log r|)^{j/2}$ on compact sets, even at collisions. Away from
pairwise distances at most $2r$, separate harmonicity and radial mollifiers
make it exactly the smooth unmollified kernel. Taking the mollification
radius to be a fixed fraction of the smallest pair distance also bounds
that kernel by a power of its logarithm. The collision neighborhood has
volume $O(r^2)$, so both integrals over it tend to zero. The mollified kernels
therefore converge locally in $L^1$ to the smooth kernel extended across
the diagonals. Their distributional limit is the original moment
distribution, proving the assertion.
\end{proof}

\begin{lemma}[Centered remainders on small revealed sets]
\label{lem:small-set-remainders}
Conditional on a stopped exploration using \cref{lem:finite-cuts}, let
$R_\delta$ be its centered flat remainders, set to zero on exposed sites.
Let $E_\delta$ be a set of sites measurable in the stopped data, in a fixed
bounded region,
and write $m_\delta$ for the area of their lattice cells. If
$m_\delta\to0$ in probability, then, for every bounded deterministic test $f$,
\[
           R_\delta(f\ind_{E_\delta})\longrightarrow0
                           \quad\text{in }L^2.
\]
The conclusion is uniform over the stopped geometries. An additional
revealed mean bounded by a fixed constant has a vanishing contribution on
these sets as well.
\end{lemma}

\begin{proof}
Condition on the stopped data and apply the second-moment case of
\cref{prop:log-moments}. If the ambient diameter is at most $R_0$, a
rearrangement bound for the logarithmic kernel gives
\[
 \E\bigl[|R_\delta(f\ind_{E_\delta})|^2\mid\mathcal E_\delta\bigr]
 \le C\|f\|_\infty^2 m_\delta^2
          \left(1+\log^+\frac{R_0^2}{m_\delta}\right),
\]
with value zero at $m_\delta=0$. To verify the discrete bound, for each
fixed site sum the decreasing logarithmic kernel first over its nearest
$O(m_\delta/\delta^2)$ cells, allowing the fixed cell-area factor of either
parity. A disk of radius comparable to
$\sqrt{m_\delta}$ bounds this sum by
$C m_\delta(1+\log^+(R_0^2/m_\delta))$; its central cell uses the cutoff
$|x-y|\vee\delta$. Sum once more over $E_\delta$. The right side is a
bounded function of $m_\delta$ tending to zero at zero, so averaging proves
the claim. A mean bounded by $M$ contributes at most
$M\|f\|_\infty m_\delta$. In particular the lemma applies to shrinking
neighborhoods of an explored set once their areas have been shown to vanish;
it does not presume that area estimate or thinness.
\end{proof}

\subsection{Harmonic moment kernels and their collision coefficients}

Fix a subsequence along which the flat fields and their moment
distributions converge, and extract the plane covariance on the same
subsequence. Write $\kappa>0$ for the coefficient denoted by $\beta$ in
\cref{prop:plane-green}:
for smooth compactly supported tests $\alpha,\beta$ of integral zero,
\begin{equation}
 \lim_{\delta\downarrow0}
 \E[K^{\mathrm{pl}}_\delta(\alpha)
       K^{\mathrm{pl}}_\delta(\beta)]
 =\kappa\iint\alpha(x)\beta(y)\log\frac1{|x-y|}\,\dd x\,\dd y.
 \label{eq:subsequential-kappa}
\end{equation}
At this stage $\kappa$ may depend on the subsequence. In particular,
neither Gaussianity nor the final normalization is being assumed.

The complementary open limit \(U\) may have several components even when
their discrete precursors belong to one component. Let \(T_j\) denote a
subsequential \(j\)th moment distribution on \(U^j\), with \(T_0=1\).
Their existence and convergence against test products follow from
\cref{lem:flat-tightness}.

We explain why separated spectator heights are legitimate exterior
observations in the first application of screening. Localize the Laplace
test to a disk \(B\), and choose \(B,Q,Q^+\) concentric, with squares
\(Q\Subset Q^+\), so that the radius of \(B\) is less than one thousandth
of the side of \(Q\).
Choose \(Q^+\) compactly inside the limiting component containing \(B\),
disjoint from the spectator disks. An arbitrary local harmonicity test can be
partitioned into tests with such small supports. Fix a simple curve surrounding \(Q\)
inside \(Q^+\setminus\overline Q\), at positive distance from both
boundaries. For all sufficiently fine meshes, choose a simple corner-grid
approximation \(\Gamma_\delta\) of this curve in the discrete chamber and
outside \(Q\).

A spectator in that chamber is read by summing corner increments from
its flat wall. Start with any such path. If it enters \(Q\), replace the
portion between its first and last encounters with \(\Gamma_\delta\) by
an arc of \(\Gamma_\delta\). The resulting path stays outside \(Q\). This
also works for spectators in a different limiting component of the same
discrete precursor. Spectators in other independently sampled chambers
may be adjoined to the exterior sigma-field: independence preserves the
screened conditional-mean bound. Thus a bounded function of all these
wall-anchored spectators is an allowed exterior observable.

On a product of disjoint interior disks, truncate the spectator product,
apply \cref{prop:screened-orthogonality}, and remove the truncation by
Holder's inequality and the higher moment bounds. This gives
\[
                     \Delta_{x_i}T_j=0\qquad(1\leq i\leq j).
\]
Summing these equations gives an elliptic equation on the product space.
Thus $T_j$ has a smooth,
separately harmonic representative off the diagonals.
\cref{lem:flat-tightness} now excludes any additional distribution on the
diagonals. Flat height
reflection gives $T_1=0$. The logarithmic moment bound and
\cref{prop:boundary-decay} apply to these representatives, by averaging
on disjoint circles. They give, respectively, polylogarithmic growth
at a collision and zero boundary trace in each variable away from the
other insertions.

\begin{lemma}[Uniform logarithmic expansion]
\label{lem:log-singularity}
If $(y,\mathbf z)$ is a collision-free list of interior points, then,
near $x=y$,
\begin{equation}
 T_j(x,y,\mathbf z)
   =L_j(y,\mathbf z)\log\frac1{|x-y|}
       +R_j(x;y,\mathbf z),
 \label{eq:logarithmic-expansion}
\end{equation}
where $R_j$ extends harmonically in $x$ across $y$. The coefficient
$L_j$ is smooth in all its arguments. On a compact set of allowed
$(y,\mathbf z)$, the expansion holds on a common radius, and the
remainder and each fixed order of its $x$-derivatives are uniformly
bounded on a smaller common radius.
\end{lemma}

\begin{proof}
A harmonic function on a punctured disk has constant angular mode
$a+b\log r$ and nonconstant modes $a_n r^n+b_n r^{-n}$, $n\geq1$.
Polylogarithmic growth forces every $b_n$ to vanish. This proves the
expansion and removability of its remainder. To obtain parameter
uniformity, choose a fixed radius $q$ avoiding all other insertions.
The coefficient is the outward flux
\begin{equation}
 L_j(y,\mathbf z)
 =-\frac1{2\pi}\int_{|w-y|=q}
                      \partial_{n_w}T_j(w,y,\mathbf z)\,\dd s_w.
 \label{eq:collision-flux}
\end{equation}
After writing $w=y+q\exp(i\vartheta)$, the integrand varies smoothly
on a collision-free compact set. This proves smoothness of $L_j$.
Subtract its logarithm from the boundary values on the circle and
apply the Poisson formula. Interior derivative estimates give the
claimed uniform bounds on $R_j$. This obtains uniformity from a
fixed circle, not from a potentially nonuniform singular limit.
\end{proof}

We next state the precise local independence needed to read this flux.
For a smooth test $\alpha$, put
$\alpha_r(x)=r^{-2}\alpha((x-y_0)/r)$. For the neutral insertions below,
use the bilinear nodal smearing of \cref{sec:spectral} in both the flat and
plane laws. Its coefficient array has exactly zero sum, and this is
precisely the discretization in \cref{eq:subsequential-kappa}.
In flat chambers \cref{lem:parity-interpolation} identifies this choice
with either parity smearing.

\begin{lemma}[A neutral pair and distant observations]
\label{lem:neutral-pair-mixing}
Fix an interior disk $B(y_0,4R)$ and a finite family of spectator
tests supported outside it. Let $Y_\delta$ be a product of the
corresponding flat height averages, or of smoothed plane increments.
If $\alpha,\beta$ are smooth zero-integral tests in disjoint bounded
disks, then
\begin{align}
 \lim_{r\downarrow0}\limsup_{\delta\downarrow0}
 \big|&\E[K_\delta(\alpha_r)K_\delta(\beta_r)Y_\delta]
 \notag\\[-1mm]
 &-\E[K^{\mathrm{pl}}_\delta(\alpha_r)
                   K^{\mathrm{pl}}_\delta(\beta_r)]\,
                      \E Y_\delta\big|=0.
 \label{eq:neutral-factorization}
\end{align}
For flat fields, the assertion is uniform over patched collections in a
common bounded region for which the local chamber contains this disk and
every spectator buffer lies in its assigned chamber. All boundary constants have been
subtracted. Spectators in other independently sampled pieces are allowed.
\end{lemma}

\begin{proof}
Here are the stopping and coloring details of the coupling. Rescale
$R$ to one and take $\rho=r^{1/2}$; all sufficiently small $r$ have
both local supports in $B(y_0,Cr)\subset B(y_0,\rho)$.
Search inward using only $(s,W)$ for a monochromatic open $W$ circuit
in the stack of annuli between radii $\rho$ and $R$. The search stops
on the first complete circuit and does not inspect its interior.
By \cref{prop:circuits,prop:comparison}, the probability of exhausting
the stack without finding a circuit is at most $C\rho^\eta$, for
some $\eta>0$. The exponent and constants depend on the fixed model
parameter and the buffered shapes, not on the flat domain.

On success reveal the entire exterior of the stopped circuit,
including its transverse colors. No $N$-cluster crosses an open
$W$ circuit, so this revelation determines the outside color
identifications without imposing an interior color constraint.
By \cref{lem:finite-cuts}, the unexplored interior is a flat law,
with its constant $s$ boundary value and independent fair colors on
its free $N$-clusters. The sign of the wall causes no change to the
law of neutral height increments: reflect a spin sign if necessary.

The spectator heights are measurable in the revealed exterior.
For a bounded flat piece, read each of them by summing corner
increments from its prescribed wall along paths avoiding
$B(y_0,R)$. Such paths exist: any path entering this interior disk
can be detoured through an annulus contained in the domain. The
same detour works in all sufficiently fine approximations, and
averaging uses only nearby exterior sites. For plane increments,
choose the integration paths between their two spectator buffers
outside the disk. Thus we have not conditioned on absolute heights
before obtaining a genuine stopping cut.

Inside the circuit apply the common-cut monotone coupling from
\cref{prop:comparison,prop:circuits} to the flat law and the plane
law. Between scales $Cr$ and $\rho$ there are order
$\log(\rho/r)$ buffered trials. Each successful common signed
$W$ circuit makes the two unexplored interiors identical in law.
After this cut, couple their $(s,W)$ configurations identically
and assign the same fair colors to the corresponding interior
$N$-clusters. Failure has probability at most $C(r/\rho)^\eta$.
This construction is valid for every exposed exterior history;
in each conditional coupling its second marginal is the same
unconditioned plane law. Consequently that marginal is independent
of the spectator observations. On a failed initial search choose
the plane copy independently. If necessary, first perform the
construction in finite flat exhaustions and then pass to the
plane law.
More explicitly, a completed $W$ cut removes every transverse interaction,
so fixing exterior colors leaves a kernel determined only by its signed
wall, with no interior color pins. For each history, bracket this kernel
and a plane exhaustion by the same ordered extreme boundary laws on an
interior buffered region. The common-cut coupling preserves both marginals;
after its cut use the identical free interior colors. Its plane marginal
does not depend on the history. Thus the circuit trials use only $(s,W)$
conditioning, never a residual law with arbitrary transverse-color pins.

The coupled local increments therefore agree outside an event
$\mathcal B_{\delta,r}$ satisfying
\begin{equation}
 \limsup_{\delta\downarrow0}\P(\mathcal B_{\delta,r})
 \leq C\{\rho^\eta+(r/\rho)^\eta\}
 \leq C r^{\eta/2}.
 \label{eq:local-coupling-error}
\end{equation}
Agreement determines local heights up to a constant, which the
exactly neutral weights remove. In particular both local smeared
variables agree on the successful event.

For completeness, the unbounded factors cost only logarithms.
For each fixed \(r\), the exactly neutral lattice coefficients are
discretizations of fixed scaled smooth profiles, with the same bounds
once \(\delta/r\) is sufficiently small. Thus
\cref{eq:flat-scaled-moments} gives, for every fixed finite $p$,
\begin{equation}
 \limsup_{\delta\downarrow0}
 \|K_\delta(\alpha_r)\|_{L^p}
 +\limsup_{\delta\downarrow0}
 \|K_\delta(\beta_r)\|_{L^p}
 \leq C_p(1+|\log r|)^{1/2}.
 \label{eq:shrinking-smearing-moment}
\end{equation}
The logarithmic diagonals are integrable, including their lattice
cells. The corresponding plane bounds are uniform in $r$, by
the scale-invariant regularity estimate
\cite[Theorem~4.5]{DKLM2026}. The fixed, smoothed spectator product
has bounded moments of every finite order. H\"older's inequality
with four factors, namely the two insertions, the spectator product,
and the failure indicator, now bounds the discrepancy by
\[
                C r^{\eta/8}(1+|\log r|)=o(1).
\]
Only original-law and plane-law moments were used. No uniform
pointwise height moments or moment bounds under an artificially
conditioned law are required. The mesh limit is always first,
with the two continuum radii fixed.
\end{proof}

\begin{proposition}[Collision coefficient]
\label{prop:collision-coefficient}
For $j\geq2$ and distinct $(y,\mathbf z)$ in $U$,
\begin{equation}
                    L_j(y,\mathbf z)=\kappa T_{j-2}(\mathbf z).
 \label{eq:collision-coefficient}
\end{equation}
\end{proposition}

\begin{proof}
Choose two smooth zero-mass tests $\alpha,\beta$ with disjoint
supports and
\[
 J=\iint\alpha(u)\beta(v)\log\frac1{|u-v|}\,\dd u\,\dd v\ne0.
\]
For an explicit choice let $\varrho_b$ be a smooth radial unit
bump supported within distance $1/10$ of $(b,0)$ and obtained by
translating one fixed bump. Take
$\alpha=\varrho_0-\varrho_1$ and
$\beta=\varrho_3-\varrho_4$. Radial harmonic averaging gives
$J=\log(8/9)\ne0$ exactly.

Take fixed radial unit mollifiers in disjoint disks around the
spectator points $\mathbf z$. By separate harmonicity, averaging
$T_{j-2}$ in those disks gives exactly $T_{j-2}(\mathbf z)$.
The same is true of $T_j$ when its first two points lie in a
small disk around $y_0$ disjoint from the spectator disks.
Apply \cref{lem:neutral-pair-mixing} and
\eqref{eq:subsequential-kappa}. For each fixed $r$, the latter
plane pairing is $\kappa J$: the added constant $\log(1/r)$
vanishes against neutral tests. Thus
\begin{equation}
 \lim_{r\downarrow0}
 \iint\alpha(u)\beta(v)
 T_j(y_0+ru,y_0+rv,\mathbf z)\,\dd u\,\dd v
       =\kappa J T_{j-2}(\mathbf z).
 \label{eq:collision-local-limit}
\end{equation}
There is no additional large-$L$ limit in this choice of tests.

On the other hand insert \cref{lem:log-singularity} in the left
side. The term
$L_j(y_0+rv,\mathbf z)\log(1/r)$ integrates to zero in $u$
\emph{exactly}, before any continuity estimate. The remaining
logarithmic term tends to $L_j(y_0,\mathbf z)J$, uniformly for
spectator configurations in compact sets. For the remainder,
subtract $R_j(y_0;y_0+rv,\mathbf z)$ inside the $u$ integral.
The subtracted term again integrates to zero, and the uniform
first-variable derivative bound gives an $O(r)$ error.
Comparison with \eqref{eq:collision-local-limit} proves the formula.

\end{proof}

\subsection{Wick's rule in the limiting chambers}

The following argument is local at
collisions and uses boundary decay, not a parametrization of the
rough boundary. It is the collision method of
\cite[Sections~9--10]{Triangular}, with the present cut and
coupling estimates supplied above.

\begin{proposition}[Gaussian moment recursion]
\label{prop:gaussian-moments}
For every $j\geq2$ and distinct points in $U$,
\begin{equation}
 T_j(x_1,\ldots,x_j)
 =\kappa\sum_{i=2}^j G_U(x_1,x_i)
          T_{j-2}(x_2,\ldots,\widehat{x_i},\ldots,x_j).
 \label{eq:bounded-wick-recursion}
\end{equation}
Consequently the limiting centered field in $U$ is a Gaussian
field with covariance $\kappa G_U$. In particular its restrictions
to different components are independent.

\end{proposition}

\begin{proof}
Fix $x_2,\ldots,x_j$ and subtract the right side of
\eqref{eq:bounded-wick-recursion}. As a function of $x_1$ the
difference is harmonic away from these points.
\cref{prop:collision-coefficient} cancels every logarithmic pole
in its component; \cref{lem:log-singularity} removes the remaining
punctures. Points in other components do not produce a pole.
The resulting harmonic function has zero boundary trace by
\cref{prop:boundary-decay}. The Dirichlet Green functions have
the same zero trace by \cref{lem:flat-green-geometry}. The difference is
bounded, since it is bounded near the boundary and its finitely
many interior punctures are removable. The maximum principle
therefore makes it zero on each component. This proves the
recursion, including mixed moments when discrete components
have separated only in the limit.

Starting with $T_0=1$ and $T_1=0$, the recursion gives zero odd
moments and the sum over perfect pairings of products
$\kappa G_U(x_i,x_j)$ for even moments. The products are locally
integrable. Since \cref{lem:flat-tightness} excludes diagonal mass,
the identity holds for the full smeared moment distributions.
These moments determine the law: if a real linear combination
$Z$ of smeared fields has variance $s^2$, then
\[
 \E Z^{2n}=(2n-1)!!s^{2n},\qquad
 \E Z^{2n+1}=0,\qquad
 \E|Z|^n\leq s^n(2n)^{n/2}.
\]
The exponential series is integrable for every real argument,
so the characteristic function is $\exp(-s^2u^2/2)$.
Apply this to all finite linear combinations and then to a
countable determining family of tests in the local Sobolev space.
The zero cross-component covariance proves independence.
\end{proof}

\subsection{Plane increment kernels and their collisions}

The free-energy correspondence used below asks for compact-uniform
convergence of all plane increment correlations. We now obtain that
conclusion from the same screened identity and neutral-pair comparison,
keeping the point correlations distinct from the smeared fields.

Use the unit corner-grid coordinates of \cref{sec:spectral}. Let
\(\widetilde K_\delta^{\mathrm{pl}}(x)\) be the piecewise constant extension
of the plane unit height used in \cite[Definition~2.4]{DKLM2026}. For an
increment \(d_i=(x_i,y_i)\), define
\[
 \Phi_{m,\delta}(d_1,\ldots,d_m)
   =\E\prod_{i=1}^m
       \bigl(\widetilde K_\delta^{\mathrm{pl}}(x_i)
             -\widetilde K_\delta^{\mathrm{pl}}(y_i)\bigr).
\]
After the fixed rotation and dilation of the grid, this is the source
correlation \(\Phi_m^{(\delta)}\) with endpoint convention
\(u_i=y_i,\ u_i'=x_i\), at the corresponding mesh. Its configuration space is
\[
 \mathcal D_m=
 \bigl\{(d_1,\ldots,d_m)\in(\R^2\times\R^2)^m:
   \{x_i,y_i\}\cap\{x_j,y_j\}=\varnothing\text{ for }i\ne j\bigr\}.
\]
The two endpoints of one increment may coincide. Put \(Q_0=1\).
Height-flip symmetry gives \(\Phi_{m,\delta}=0\) for odd \(m\), so their
limits \(Q_m\) are zero. For even \(m\), the regularity estimate and
precompactness result \cite[Theorem~4.5 and Corollary~4.6]{DKLM2026} give,
after a diagonal extraction on the chosen subsequence,
\[
              \Phi_{m,\delta}\longrightarrow Q_m
              \quad\text{uniformly on compact subsets of }\mathcal D_m
\]
at every order. Each \(Q_m\) is continuous there. These are correlation
kernels, not products of point values of a continuum field.

We first relate these kernels to the nodal smooth insertions used above.
For each neutral smooth test \(\alpha_i\), choose a smooth mass-one density
\(\rho_i\) in a separate remote disk, with all these disks disjoint from
the test supports and from each other. Let \(b_{\delta,v}^{(i)}\) be the
nodal interpolation coefficients for \(\alpha_i\), and let
\(r_{\delta,w}^{(i)}\) be discrete weights approximating \(\rho_i\) with
sum one. The corresponding signed measures converge to
\(\alpha_i(x)\dd x\) and \(\rho_i(y)\dd y\), and their atoms have absolute
weights \(O(\delta^2)\). Exact neutrality gives the finite identity
\[
 \sum_v b_{\delta,v}^{(i)}K(v)
 =\sum_{v,w} b_{\delta,v}^{(i)}r_{\delta,w}^{(i)}
                    \bigl(K(v)-K(w)\bigr).
\]
Thus products of nodal insertions test \(\Phi_{m,\delta}\) with
independently integrated reference variables, rather than coincident
fixed anchors. Away from collisions between endpoints of different
increments, use compact-uniform convergence. The regularity estimate
bounds the omitted neighborhoods by integrable powers of logarithms,
uniformly after the lattice cutoff; their exact coincidence cells
contribute at most \(C\delta^2(1+|\log\delta|)^q\) for a fixed \(q\).
These integrals vanish as the neighborhoods shrink. Coincidence of an
increment's own endpoints has the vanishing increment-collapse bound,
not a singular term. This identifies the full smeared moment distributions
with the integrals of \(Q_m\), without losing a contribution at an anchor
or a collision.

For \(m=2\), this identification and \cref{eq:subsequential-kappa} give
the covariance kernel
\begin{equation}
 C_\kappa(d_i,d_j)
   =\kappa\log
       \frac{|x_i-y_j|\,|y_i-x_j|}
            {|x_i-x_j|\,|y_i-y_j|}.
 \label{eq:increment-covariance}
\end{equation}
as \(Q_2\). Indeed one may approximate each endpoint by a smooth mass-one
bump in a disjoint disk; the finite increment additivity and continuity of
\(Q_2\) then recover the displayed point kernel from the neutral pairings.

Screening gives separate harmonicity of \(Q_m\) in every endpoint away
from the endpoints of other increments. To see the required exterior
measurability, localize a Laplace test in \(B\Subset Q\Subset Q^+\) as
before, with \(Q^+\) disjoint from the other endpoint buffers. In the
differentiated increment the remote endpoint cancels against the exactly
neutral Laplace coefficients. Each remaining increment is read along a
path between its two buffers outside \(Q\), using a fixed annular detour
when necessary. Smooth the remote endpoints in their buffers, apply
\cref{prop:screened-orthogonality} to the truncated spectator product, and
remove the truncation by the regularity moment bounds. The smearing
identification just proved gives distributional harmonicity of \(Q_m\).
Elliptic regularity gives its smooth separately harmonic representative
where endpoints of different increments are distinct. The argument also
allows the two endpoints of the differentiated increment to coincide,
since its reference term has canceled.

Near a collision between different increments, the source regularity
estimate bounds \(Q_m\) by a power of the logarithm of the separation.
The analytic proof of \cref{lem:log-singularity} therefore applies to
each such endpoint collision: there is one logarithmic term, its coefficient
is given by a flux on a fixed circle, and the harmonic remainder and its
derivatives are uniformly bounded on compact sets of the other endpoints.
We can now identify that coefficient.

\begin{lemma}[Signed plane collision coefficients]
\label{lem:plane-collision-coefficient}
Fix \(m\ge2\), and keep all endpoints other than \(x_1\) distinct. For
\(i\ge2\), the coefficient of \(\log(1/|x_1-z|)\) in \(Q_m\) is
\[
 \begin{cases}
  \kappa Q_{m-2}(d_2,\ldots,\widehat{d_i},\ldots,d_m),&z=x_i,\\
  -\kappa Q_{m-2}(d_2,\ldots,\widehat{d_i},\ldots,d_m),&z=y_i.
 \end{cases}
\]
After subtracting this logarithm, the remainder extends harmonically
through the collision. Swapping \(x_1\) and \(y_1\) gives the corresponding
rule with the opposite sign for the negative endpoint.
\end{lemma}

\begin{proof}
First consider two positive endpoints. Use the neutral tests
\(\alpha,\beta\) from the proof of \cref{prop:collision-coefficient}, whose
logarithmic pairing is \(J=\log(8/9)\). Smooth every remote endpoint in a
small fixed radial disk. Integrating the two colliding endpoints against
\(\alpha_r,\beta_r\) cancels their two remote reference terms exactly.
By the smearing identification, the resulting expression is the limit
of the two neutral lattice insertions times a product of smoothed plane
increments. These spectators have the bounded moments required by
\cref{lem:neutral-pair-mixing}. That lemma and
\cref{eq:subsequential-kappa} give the limit
\[
       \kappa J\,Q_{m-2}(d_2,\ldots,\widehat{d_i},\ldots,d_m)
       \qquad(r\downarrow0).
\]
Separate harmonicity removes the spectator mollifiers. In the uniform
logarithmic expansion, the term containing \(\log(1/r)\) integrates to
zero against \(\alpha\) before any limit; smoothness of the coefficient
and the uniform first-variable derivative bound for the remainder leave
the coefficient times \(J\). Comparing the two expressions proves the
positive sign. Antisymmetry under reversal of either increment gives the
negative signs. This calculation uses only smoothed spectator moments.
\end{proof}

\begin{proposition}[Plane Gaussian correlations]
\label{prop:plane-gaussian-moments}
Along the chosen subsequence, all plane increment correlation functions
satisfy Wick's rule with covariance \(C_\kappa\). More precisely,
\(\Phi_{m,\delta}\to Q_m\) uniformly on compact subsets of \(\mathcal D_m\)
at every order, and \(Q_m\) is zero for odd \(m\), while for even \(m\)
it is the sum over pairings of products of \(C_\kappa(d_i,d_j)\).
\end{proposition}

\begin{proof}
The cases \(m=0,1,2\) have been established. For even \(m\ge4\), induct
on \(m\) and subtract
\begin{equation}
 \sum_{i=2}^m C_\kappa(d_1,d_i)
                 Q_{m-2}(d_2,\ldots,\widehat{d_i},\ldots,d_m)
 \label{eq:plane-wick-recursion}
\end{equation}
from $Q_m$. In the variable $x_1$,
\cref{lem:plane-collision-coefficient} cancels all logarithmic poles at
$x_i$ and $y_i$, $i\geq2$, and its harmonic remainders remove the
punctures. The regularity estimate gives continuous extension through
$x_1=y_1$, with
value zero there; this is the collapse of the first increment,
not an additional collision between different factors. The
remainder is consequently entire harmonic. With all other
endpoints fixed, the same estimate bounds its growth by a fixed
power of $1+\log(2+|x_1|)$, in particular by $o(|x_1|)$.
An entire harmonic function of sublinear growth is constant:
the interior gradient bound on disks of radius $M$, followed
by $M\to\infty$, makes its gradient zero. At $x_1=y_1$ both
$Q_m$ and \eqref{eq:plane-wick-recursion} vanish. The constant
is therefore zero. This proves the plane Wick recursion.

The argument initially identifies the kernels where all the
endpoints are distinct. The quoted equicontinuity and the
increment-collapse estimate extend the identity locally
uniformly to the full configuration spaces allowing the two
endpoints of one increment to coincide. Thus all orders, not
only covariance or smooth-test laws, have been identified on
one common subsequence.
\end{proof}

\subsection{Identifying the coefficient without circularity}

\begin{proposition}[Plane normalization]
\label{prop:plane-normalization}
For every subsequence considered above,
\begin{equation}
 \kappa=\frac1{2\pi\arcsin(c/2)}=\frac1{a^2}.
 \label{eq:identified-kappa}
\end{equation}
\end{proposition}

\begin{proof}
The coefficient is strictly positive by \cref{prop:plane-green}.

Set $\sigma^2=2\pi\kappa$. The plane Wick kernels just proved
are precisely those of $\sigma$ times a plane GFF in the
normalization with singularity $-(2\pi)^{-1}\log|x-y|$.
The fixed rotation and dilation used above leave the increment
cross-ratio in \cref{eq:increment-covariance} unchanged.
Their compact-uniform convergence at every order is the exact
hypothesis of \cite[Theorem~8.1]{DKLM2026}. That convergence
criterion, whose input is the regularity estimate, supplies
the full-plane subsequential scaling limit in the sense required
by its Definition~2.7. We may therefore apply
\cite[Theorems~4.3 and~4.4]{DKLM2026}: their free-energy correspondence
and explicit computation give
\[
                    \sigma^2=\frac1{\arcsin(c/2)}.
\]
The source uses the column slope $I_{\rm col}/L$ and per-site free
energy, as reconciled in \cref{sec:input-amplitude}; that appendix
also records the finite-error factors before their final limits.
All three cited inputs apply for $1\leq c\leq2$; the present
fixed parameter satisfies $1<c\leq2$. This use does not invoke
the narrower-range height convergence theorem of that paper.
It follows that $2\pi\kappa=1/\arcsin(c/2)$. Finally
$a^2=g(\pi/2)^2=2\pi\arcsin(c/2)$, proving
\eqref{eq:identified-kappa}.

The logical order is important. The spectral argument supplied
only a subsequential coefficient; the local collision argument
identified every Gaussian correlation with that coefficient;
the convergence criterion established the requisite plane GFF
limit; only now has the free energy fixed the coefficient.
\end{proof}

\subsection{Completion of the flat limit and conditioning at a cut}

\begin{proof}[Proof of \cref{thm:flat-gff}]
By \cref{lem:flat-green-geometry}, every compact subset of $U$ eventually
lies inside the discrete open set.
On a finite list of interior buffers, \cref{lem:flat-tightness} gives joint
subsequential field and moment convergence. The screened identity applies
there with their fixed positive clearances. The boundary estimate uses the
connected accessible wall in each chamber, and remains valid at every
boundary point of the limiting components, as proved in
\cref{prop:boundary-decay}. The Green boundary estimate in
\cref{lem:flat-green-geometry} supplies the matching zero trace.
Thus the proofs of
\cref{prop:collision-coefficient,prop:gaussian-moments} apply to these
components without a smooth-boundary assumption. They identify all mixed
moment kernels, including those in different components, with the block
Dirichlet covariance $\kappa G_U$.

The plane coefficient on any simultaneous subsequence is
$\kappa=1/a^2$ by \cref{prop:plane-normalization}. Hence all subsequential
laws agree. Tightness proves full convergence, and the moment bounds give
convergence of moments. \cref{lem:parity-interpolation} supplies the asserted
parity and interpolation equivalences. For a mesh-dependent number of
independent chambers, the moment proof is unchanged: each retained site
belongs to at most one chamber, and shared-cluster counts between different
chambers are zero.
In particular the constants do not accumulate with the number of chambers.
The same uniform bounds, with zero extension, prove the final tightness
assertion. That assertion does not identify an extra field supported on a
possibly positive-area limiting exposed set; subsequent uses of thinness
will address the full decomposition across that set.
\end{proof}

At a stopped cut, let \(A_\delta\) be the revealed closed set with its
lattice-scale thickening. On any fixed compact test region \(Q\) separated
from the initial domain boundary, the closed complement of the resolved
chamber faces satisfies
\[
 F_\delta\cap Q\subset A_\delta^{C\delta},
 \qquad
 A_\delta\cap Q\subset F_\delta^{C\delta},
\]
where \(E^\eta=\{x:\dist(x,E)<\eta\}\). Both descriptions contain the same
exposed cells and perimeter edges; only the one-cell medial displacement
differs. These inclusions transfer the neighborhood estimates in
\cref{sec:local-sets} between the revealed cut and the resolved complement.

\begin{proposition}[Conditional flat remainders]
\label{prop:conditional-flat-gff}
Suppose a stopping exploration uses the cuts in \cref{lem:finite-cuts}, so
that, conditional on its revealed data $\mathcal E_\delta$, the unexplored
chambers have independent flat laws with pairwise disjoint interiors. Let
$R_\delta$ be their centered unit-height remainders, patched by zero, and
let $U_\delta,F_\delta$ be their open union and closed complement in
$\overline B_0$. Fix real tests $f_1,\ldots,f_\ell\in C_c^\infty(B_0)$,
$u_1,\ldots,u_\ell\in\R$, and $d_0>0$. Put
\[
 A_{\delta,d_0}
 =\left\{\dist\left(\bigcup_i\supp f_i,F_\delta\right)\ge d_0\right\}.
\]
Then
\begin{equation}
\begin{split}
 \E\Bigg[\ind_{A_{\delta,d_0}}
 \Bigg|&
 \E\left[\exp\left(i\sum_i u_iR_\delta(f_i)\right)
                  \,\middle|\,\mathcal E_\delta\right]\\
 &-\exp\left(-\frac1{2a^2}
                  \sum_{i,j}u_i u_jG_{U_\delta}(f_i,f_j)\right)
 \Bigg|\Bigg]\longrightarrow0.
\end{split}
\label{eq:conditional-flat-characteristic}
\end{equation}
Here \(G_{U_\delta}\) is the continuum Dirichlet kernel of the resolved
open chamber union, with zero covariance between its components.

Multiplying the displayed difference by any uniformly bounded
\(\mathcal E_\delta\)-measurable variable preserves the same estimate with
its clearance indicator. To state the limiting consequence, let
\(\mathsf E\) be a Polish space and let \(e_\delta\) be an
\(\mathcal E_\delta\)-measurable \(\mathsf E\)-valued variable. Suppose,
for one fixed \(s<-1\), that along a subsequence
\[
 (e_\delta,F_\delta,R_\delta)\Longrightarrow(e,F,R)
 \quad\text{in }\mathsf E\times\mathcal K(\overline B_0)
                   \times H^s_{\mathrm{loc}}(B_0),
\]
where \(\mathcal K(\overline B_0)\) is the space of nonempty compact sets
with the Hausdorff metric. The limit \(F\) is almost surely connected and
contains \(\partial B_0\). Put \(U=B_0\setminus F\) and
\[
 A_F=\left\{\dist\left(\bigcup_i\supp f_i,F\right)>0\right\}.
\]
Then, almost surely,
\begin{equation}
\begin{split}
 &\ind_{A_F}\,
 \E\left[\exp\left(i\sum_i u_iR(f_i)\right)
             \,\middle|\,\sigma(e,F)\right]\\
 &\qquad=\ind_{A_F}\,
 \exp\left(-\frac1{2a^2}
                \sum_{i,j}u_i u_jG_U(f_i,f_j)\right).
\end{split}
\label{eq:conditional-flat-limit}
\end{equation}
Thus, conditional on \(\sigma(e,F)\), the restrictions of \(R\) to the
component interiors of \(U\) are independent zero-boundary GFFs divided by
\(a\), including components born from a pinched single precursor. This is
first tested on deterministic compact buffers with positive limiting
clearance and then extended by a countable exhaustion. The retained variable
\(e_\delta\) may encode any exact-cut coordinates and flat constants for
which this joint convergence holds; the constants are added after the
centered assertion. Its \(\mathcal E_\delta\)-measurability excludes data
from any later revelation.
\end{proposition}

\begin{proof}
For a deterministic realized chamber configuration, let its flat product
law supply the characteristic function in the first line of
\cref{eq:conditional-flat-characteristic}. The error from the displayed
Gaussian characteristic function tends to zero uniformly over all such
configurations with the stated clearance. Otherwise choose a sequence with
error bounded away from zero. Compactness of the closed complements in
\(\overline B_0\) gives a Hausdorff-convergent further subsequence, and the
clearance persists in its limit. The deterministic theorem gives convergence
of the flat characteristic function to the Gaussian one for that limit.
\Cref{lem:flat-green-geometry} gives
\(G_{U_\delta}(f_i,f_j)\to G_U(f_i,f_j)\) on the same subsequence.
The two characteristic functions therefore have the same limit, a
contradiction.

Condition now on the stopped data. By \cref{lem:finite-cuts}, the
conditional kernel is exactly the deterministic product kernel just
considered, with its constants subtracted. The uniform error proves
\cref{eq:conditional-flat-characteristic}. Multiplication by any bounded
stopped-measurable variable preserves convergence to zero.

Assume the displayed joint convergence and write
\(d(F)=\dist(\bigcup_i\supp f_i,F)\). Choose a bounded continuous function
\(\chi\) of the clearance whose support is contained in \([d_0,\infty)\)
for some \(d_0>0\), and a bounded continuous function
\(v\) on \(\mathsf E\times\mathcal K(\overline B_0)\).
Multiply the conditional difference by
\(\chi(d(F_\delta))v(e_\delta,F_\delta)\), and use the tower property for
the term involving \(R_\delta\). The estimate just proved makes the
resulting expectation tend to zero. Evaluation of \(R\) on a smooth test
and the clearance \(d(F)\) are continuous in the displayed state space.
The connected complements containing \(\partial B_0\) form a closed
subspace supporting every \(F_\delta\) and \(F\). On that subspace,
\cref{lem:flat-green-geometry} makes the Gaussian expression continuous
where \(\chi\) is supported. Joint convergence therefore gives
\begin{align*}
 &\E\left[\chi(d(F))v(e,F)
       \exp\left(i\sum_i u_iR(f_i)\right)\right]\\
 &\qquad=\E\left[\chi(d(F))v(e,F)
       \exp\left(-\frac1{2a^2}
                 \sum_{i,j}u_i u_jG_U(f_i,f_j)\right)\right].
\end{align*}
The monotone class theorem on the Polish space
\(\mathsf E\times\mathcal K(\overline B_0)\) extends this identity to
bounded measurable \(v\). A countable family of cutoffs \(\chi\) whose
positive sets cover \(d(F)>0\) now gives
\cref{eq:conditional-flat-limit}. A countable determining family of
smooth tests then identifies the regular conditional law on the component
interiors. The zero cross-component Green kernel gives their independence,
including when one precursor separates into several limiting components.
Increasing the list of buffers and decreasing their required clearance
gives the full interior statement. Jointly retained constants can then be
added.
\end{proof}
